\documentclass[CJK]{beamer}
\usepackage{CJKutf8}
\usepackage{beamerthemesplit}
\usetheme{Malmoe}
\useoutertheme[footline=authortitle]{miniframes}
\usepackage{amsmath}
\usepackage{amssymb}
\usepackage{graphicx}
\usepackage{eufrak}
\usepackage{color}
\usepackage{slashed}
\usepackage{simplewick}
\usepackage{tikz}
\usepackage{tcolorbox}
\graphicspath{{../figures/}}
%%figures
\def\lfig#1#2{\includegraphics[width=#1 in]{#2}}
\def\tfig#1#2{\includegraphics[height=#1 in]{#2}}
\def\addfig#1#2{\begin{center}\includegraphics[width=#1 in]{#2}\end{center}}
\def\wulian{\includegraphics[width=0.18in]{emoji_wulian.jpg}}
\def\bigwulian{\includegraphics[width=0.35in]{emoji_wulian.jpg}}
\def\bye{\includegraphics[width=0.18in]{emoji_bye.jpg}}
\def\bigbye{\includegraphics[width=0.35in]{emoji_bye.jpg}}
\def\huaixiao{\includegraphics[width=0.18in]{emoji_huaixiao.jpg}}
\def\bighuaixiao{\includegraphics[width=0.35in]{emoji_huaixiao.jpg}}
\def\jianxiao{\includegraphics[width=0.18in]{emoji_jianxiao.jpg}}
\def\bigjianxiao{\includegraphics[width=0.35in]{emoji_jianxiao.jpg}}
%% colors
\def\blacktext#1{{\color{black}#1}}
\def\bluetext#1{{\color{blue}#1}}
\def\redtext#1{{\color{red}#1}}
\def\darkbluetext#1{{\color[rgb]{0,0.2,0.6}#1}}
\def\skybluetext#1{{\color[rgb]{0.2,0.7,1.}#1}}
\def\cyantext#1{{\color[rgb]{0.,0.5,0.5}#1}}
\def\greentext#1{{\color[rgb]{0,0.7,0.1}#1}}
\def\darkgray{\color[rgb]{0.2,0.2,0.2}}
\def\lightgray{\color[rgb]{0.6,0.6,0.6}}
\def\gray{\color[rgb]{0.4,0.4,0.4}}
\def\blue{\color{blue}}
\def\red{\color{red}}
\def\green{\color{green}}
\def\darkgreen{\color[rgb]{0,0.4,0.1}}
\def\darkblue{\color[rgb]{0,0.2,0.6}}
\def\skyblue{\color[rgb]{0.2,0.7,1.}}
%%control
\def\be{\begin{equation}}
\def\ee{\nonumber\end{equation}}
\def\bea{\begin{eqnarray}}
\def\eea{\nonumber\end{eqnarray}}
\def\bch{\begin{CJK}{UTF8}{gbsn}}
\def\ech{\end{CJK}}
\def\bitem{\begin{itemize}}
\def\eitem{\end{itemize}}
\def\bcenter{\begin{center}}
\def\ecenter{\end{center}}
\def\bex{\begin{minipage}{0.2\textwidth}\includegraphics[width=0.6in]{jugelizi.png}\end{minipage}\begin{minipage}{0.76\textwidth}}
\def\eex{\end{minipage}}
\def\chtitle#1{\frametitle{\bch#1\ech}}
\def\bmat#1{\left(\begin{array}{#1}}
\def\emat{\end{array}\right)}
\def\bcase#1{\left\{\begin{array}{#1}}
\def\ecase{\end{array}\right.}
\def\bmini#1{\begin{minipage}{#1\textwidth}}
\def\emini{\end{minipage}}
\def\tbox#1{\begin{tcolorbox}#1\end{tcolorbox}}
\def\pfrac#1#2#3{\left(\frac{\partial #1}{\partial #2}\right)_{#3}}
%%symbols
\def\newt#1#2{\left(\begin{array}{c}#1\\ #2\end{array}\right)}
\def\reof#1{\,\mathrm{Re}{\left(#1\right)\,}}
\def\imof#1{\,\mathrm{Im}{\left(#1\right)\,}}
\def\Arg#1{\,\mathrm{Arg}\,#1\,}
\def\bropt{\,(\ \ \ )}
\def\sone{$\star$}
\def\stwo{$\star\star$}
\def\sthree{$\star\star\star$}
\def\sfour{$\star\star\star\star$}
\def\sfive{$\star\star\star\star\star$}
\def\rint{{\int_\leftrightarrow}}
\def\roint{{\oint_\leftrightarrow}}
\def\stdHf{{\textit{\r H}_f}}
\def\deltaH{{\Delta \textit{\r H}}}
\def\ii{{\dot{\imath}}}
\def\skipline{{\vskip0.1in}}
\def\skiplines{{\vskip0.2in}}
\def\lagr{{\mathcal{L}}}
\def\hamil{{\mathcal{H}}}
\def\vecv{{\mathbf{v}}}
\def\vecx{{\mathbf{x}}}
\def\vecy{{\mathbf{y}}}
\def\veck{{\mathbf{k}}}
\def\vecp{{\mathbf{p}}}
\def\vecn{{\mathbf{n}}}
\def\vecA{{\mathbf{A}}}
\def\vecP{{\mathbf{P}}}
\def\vecsigma{{\mathbf{\sigma}}}
\def\hatJn{{\hat{J_\vecn}}}
\def\hatJx{{\hat{J_x}}}
\def\hatJy{{\hat{J_y}}}
\def\hatJz{{\hat{J_z}}}
\def\hatj#1{\hat{J_{#1}}}
\def\hatphi{{\hat{\phi}}}
\def\hatq{{\hat{q}}}
\def\hatpi{{\hat{\pi}}}
\def\vel{\upsilon}
\def\Dint{{\mathcal{D}}}
\def\adag{{\hat{a}^\dagger}}
\def\bdag{{\hat{b}^\dagger}}
\def\cdag{{\hat{c}^\dagger}}
\def\ddag{{\hat{d}^\dagger}}
\def\hata{{\hat{a}}}
\def\hatb{{\hat{b}}}
\def\hatc{{\hat{c}}}
\def\hatd{{\hat{d}}}
\def\hatN{{\hat{N}}}
\def\hatH{{\hat{H}}}
\def\hatp{{\hat{p}}}
\def\Fup{{F^{\mu\nu}}}
\def\Fdown{{F_{\mu\nu}}}
\def\newl{\nonumber \\}
\def\vece{\mathrm{e}}
\def\calM{{\mathcal{M}}}
\def\calT{{\mathcal{T}}}
\def\calR{{\mathcal{R}}}
\def\barpsi{\bar{\psi}}
\def\baru{\bar{u}}
\def\barv{\bar{\upsilon}}
\def\qeq{\stackrel{?}{=}}
\def\torder#1{\mathcal{T}\left(#1\right)}
\def\rorder#1{\mathcal{R}\left(#1\right)}
\def\contr#1#2{\contraction{}{#1}{}{#2}#1#2}
\def\trof#1{\mathrm{Tr}\left(#1\right)}
\def\trace{\mathrm{Tr}}
\def\comm#1{\ \ \ \left(\mathrm{used}\ #1\right)}
\def\tcomm#1{\ \ \ (\text{#1})}
\def\slp{\slashed{p}}
\def\slk{\slashed{k}}
\def\calp{{\mathfrak{p}}}
\def\veccalp{\mathbf{\mathfrak{p}}}
\def\Tthree{T_{\tiny \textcircled{3}}}
\def\pthree{p_{\tiny \textcircled{3}}}
\def\dbar{{\,\mathchar'26\mkern-12mu d}}
\def\erf{\mathrm{erf}}
\def\const{\mathrm{constant}}
\def\pheat{\pfrac p{\ln T}V}
\def\vheat{\pfrac V{\ln T}p}
%%units
\def\fdeg{{^\circ \mathrm{F}}}
\def\cdeg{^\circ \mathrm{C}}
\def\atm{\,\mathrm{atm}}
\def\angstrom{\,\text{\AA}}
\def\SIL{\,\mathrm{L}}
\def\SIkm{\,\mathrm{km}}
\def\SIyr{\,\mathrm{yr}}
\def\SIGyr{\,\mathrm{Gyr}}
\def\SIV{\,\mathrm{V}}
\def\SImV{\,\mathrm{mV}}
\def\SIeV{\,\mathrm{eV}}
\def\SIkeV{\,\mathrm{keV}}
\def\SIMeV{\,\mathrm{MeV}}
\def\SIGeV{\,\mathrm{GeV}}
\def\SIcal{\,\mathrm{cal}}
\def\SIkcal{\,\mathrm{kcal}}
\def\SImol{\,\mathrm{mol}}
\def\SIN{\,\mathrm{N}}
\def\SIHz{\,\mathrm{Hz}}
\def\SIm{\,\mathrm{m}}
\def\SIcm{\,\mathrm{cm}}
\def\SIfm{\,\mathrm{fm}}
\def\SImm{\,\mathrm{mm}}
\def\SInm{\,\mathrm{nm}}
\def\SImum{\,\mathrm{\mu m}}
\def\SIJ{\,\mathrm{J}}
\def\SIW{\,\mathrm{W}}
\def\SIkJ{\,\mathrm{kJ}}
\def\SIs{\,\mathrm{s}}
\def\SIkg{\,\mathrm{kg}}
\def\SIg{\,\mathrm{g}}
\def\SIK{\,\mathrm{K}}
\def\SImmHg{\,\mathrm{mmHg}}
\def\SIPa{\,\mathrm{Pa}}
%page
\def\secpage#1#2{\begin{frame}\bch\bcenter{\bf \Huge #1} \skipline \tbox{#2}\ecenter\ech\end{frame}}

\def\courseurl{https://github.com/zqhuang/SYSU\_TD}

\def\tpage#1#2{
\begin{frame}
\begin{center}
\begin{Large}
\bch
热学 \\
第#1讲 #2

{\vskip 0.3in}

黄志琦

\ech
\end{Large}
\end{center}

\vskip 0.2in

\bch
教材：《热学》第二版，赵凯华，罗蔚茵，高等教育出版社
\ech

\bch
课件下载
\ech
\courseurl
\end{frame}
}

\def\bfr#1{
\begin{frame}
\chtitle{#1} 
\bch
}

\def\efr{
\ech 
\end{frame}
}

  \date{}
\begin{document}
\tpage{17}{Bessel Functions of the First Kind II}

\begin{frame}
\chtitle{本讲内容}
\bch
\bitem
\item{圆形薄膜振动问题的一般初始条件}
\item{圆盘上的热传导问题}
\item{无边界的圆对称问题}      
\eitem
\ech
\end{frame}


\section{Wave Equation with general IC}

\begin{frame}
  \chtitle{圆形薄膜振动问题的一般初始条件}
  \bch
  \addfig{1.}{drum.jpg}
  考虑边界固定，半径为$R$的圆形薄膜的横向小振动问题。在$t=0$时刻的初始位移为$f(r,\theta)$ ，初始速度为$g(r,\theta)$。其中$(r,\theta)$是以圆盘中心为原点建立的极坐标。$f$和$g$都满足边界条件：$f(r,\theta+2\pi) = f(r,\theta)$，$g(r,\theta+2\pi)=g(r,\theta)$，$f(R,\theta) = g(R,\theta) = 0$。求解之后薄膜的振动。
    \ech
\end{frame}

\begin{frame}
  \chtitle{写出方程和边界条件}
  \bch
  设位移为$u(r,\theta,t)$。

  \bea
  \frac{\partial^2u}{\partial t^2} - a^2\nabla^2 u = 0 , \newl
  \left.u\right\vert_{r=R} = 0,\newl
  \left.u\right\vert_{t = 0} = f(r,\theta) , \newl
  \left.\frac{\partial u}{\partial t}\right\vert_{t = 0} = g(r,\theta) .
  \eea
  \ech
\end{frame}


\begin{frame}
  \chtitle{分离变量}
  \bch
  符合边界条件的分离变量形式解为：
  $$ J_m\left(\frac{\mu_{m,i}r}{R}\right)\cos{(m\theta)} \cos{\left(\frac{\mu_{m,i}at}{R}\right)},\ J_m\left(\frac{\mu_{m,i}r}{R}\right)\cos{(m\theta)} \sin{\left(\frac{\mu_{m,i}at}{R}\right)}, $$
  其中$m\in Z, \ \mu_{m,i}$是$J_m(x)$的第$i$个正实数根。 当$m>0$时还有
  $$ J_m\left(\frac{\mu_{m,i}r}{R}\right)\sin{(m\theta)} \cos{\left(\frac{\mu_{m,i}at}{R}\right)},\ J_m\left(\frac{\mu_{m,i}r}{R}\right)\sin{(m\theta)} \sin{\left(\frac{\mu_{m,i}at}{R}\right)}. $$
  
  \ech
\end{frame}


\begin{frame}
  \chtitle{级数展开}
  \bch
  完整的级数形式解为
  \bea
  u &=& \sum_{i=1}^\infty \left\{\frac{A_{0,i}}{2} J_0\left(\frac{\mu_{0,i}r}{R}\right) \cos{\left(\frac{\mu_{0,i}at}{R}\right)} + \frac{B_{0,i}}{2} J_0\left(\frac{\mu_{0,i}r}{R}\right) \sin{\left(\frac{\mu_{0,i}at}{R}\right)}\right. \newl
  && +\sum_{m=1}^\infty \left[ A_{m,i} J_m\left(\frac{\mu_{m,i}r}{R}\right)\cos{(m\theta)} \cos{\left(\frac{\mu_{m,i}at}{R}\right)} \right. \newl
  && + B_{m,i} J_m\left(\frac{\mu_{m,i}r}{R}\right)\cos{(m\theta)} \sin{\left(\frac{\mu_{m,i}at}{R}\right)} \newl
    && + C_{m,i} J_m\left(\frac{\mu_{m,i}r}{R}\right)\sin{(m\theta)} \cos{\left(\frac{\mu_{m,i}at}{R}\right)} \newl
      &&  +\left.\left. D_{m,i} J_m\left(\frac{\mu_{m,i}r}{R}\right)\sin{(m\theta)} \sin{\left(\frac{\mu_{m,i}at}{R}\right)} \right]\right\}
  \eea
  (在$A_{0,i}$和$B_{0,i}$里的$\frac{1}{2}$因子是为了和$m\ne 0$的结果看起来更一致。)
  \ech
\end{frame}



\begin{frame}
  \chtitle{初始位移条件}
  \bch
  利用初始位移条件得到
  {\small
  \bea
  f(Rx,\theta)  &=& \sum_{i=1}^\infty \frac{A_{0,i}}{2} J_0\left(\mu_{0,i}x\right)  \newl
  &+& \sum_{m=1}^\infty\sum_{i=1}^\infty \left[ A_{m,i} J_m\left(\mu_{m,i}x\right)\cos{(m\theta)} + C_{m,i} J_m\left(\mu_{m,i}x\right)\sin{(m\theta)}  \right]
  \eea }
  其中$x = \frac{r}{R}$。

  由谐函数的正交性得到
  $$ A_{m,i} = \frac{2}{\pi\left[J_{m+1}(\mu_{m,i})\right]^2} \int_0^1 xJ_m(\mu_{m,i} x) dx\int_0^{2\pi} \cos{(m\theta)}   f(Rx, \theta) d\theta $$
  $$C_{m,i} = \frac{2}{\pi\left[J_{m+1}(\mu_{m,i})\right]^2} \int_0^1 xJ_m(\mu_{m,i} x) dx\int_0^{2\pi} \sin{(m\theta)}   f(Rx, \theta) d\theta $$
  \ech
\end{frame}


\begin{frame}
  \chtitle{初始速度条件}
  \bch
  利用初始速度条件得到
  {\small
  \bea
  g(Rx,\theta)  &=& \sum_{i=1}^\infty \frac{B_{0,i}\mu_{0,i}a}{2R} J_0\left(\mu_{0,i}x\right)  + \sum_{m=1}^\infty\sum_{i=1}^\infty \left[ \frac{B_{m,i}\mu_{m,i}a}{R} J_m\left(\mu_{m,i}x\right)\cos{(m\theta)} \right. \newl
    && \left. + \frac{D_{m,i}\mu_{m,i}a}{R} J_m\left(\mu_{m,i}x\right)\sin{(m\theta)}  \right]
  \eea }
  其中$x = \frac{r}{R}$。

  由谐函数的正交性得到
  $$ B_{m,i} = \frac{2R}{\pi \mu_{m,i}a\left[J_{m+1}(\mu_{m,i})\right]^2} \int_0^1 xJ_m(\mu_{m,i} x) dx\int_0^{2\pi} \cos{(m\theta)}   g(Rx, \theta) d\theta $$
  $$D_{m,i} = \frac{2R}{\pi\mu_{m,i}a\left[J_{m+1}(\mu_{m,i})\right]^2} \int_0^1 xJ_m(\mu_{m,i} x) dx\int_0^{2\pi} \sin{(m\theta)}   g(Rx, \theta) d\theta $$
  \ech
\end{frame}



\section{Heat Equation on a Disk}

\begin{frame}
  \chtitle{}
  \bch

  有孤立的，半径为$R$的均匀厚度圆形金属片。初态$t=0$时刻金属片中心温度为$T_0$，金属片上距离中心$r$处的温度为
$$T(r) =T_0\left(1+\frac{r^2}{R^2}\right).$$
已知金属片质量为$\rho$，导热系数为$\lambda$，单位质量比热为$c$。求之后圆盘上各点的温度变化。

  \ech
\end{frame}


\begin{frame}
  \chtitle{写出方程}
  \bch
\bea
\frac{\partial^2T}{\partial t^2}-a\nabla^2 T &=& 0, \newl
\left.\frac{\partial T}{\partial r}\right\vert_{r=a} &=& 0, \newl
\left.T\right\vert_{t=0} &=& T_0\left(1+\frac{r^2}{R^2}\right).
\eea
  \ech
\end{frame}



\section{Homework}

\begin{frame}
  \chtitle{课后作业(题号38-40)}
  \bch
  \bitem
\item[35]{通过逐项求导的办法证明贝塞尔函数的两个递推关系}
\item[36]{如果$J_m(\mu) = 0$ （$\mu$为正实数，$m$为非负整数)，计算积分
    $$\int_0^1x\left[J_{m+1}(\mu x)\right]^2 dx. $$
}
\item[37]{ 考虑半径为$R$的圆形薄膜的横向小振动问题。在$t=0$时刻的初始速度为$A\left[1-\left(\frac{r}{R}\right)^2\right]$ （其中$A$为常量，$r$为距离圆盘中心的距离，$0\le r\le R$)，初始位移为零。求解之后薄膜的振动。
 }
  \eitem
  \ech
\end{frame}


\end{document}
